\section{A discriminating experiment on attention and report}
The following design is prospective. No participant data or effect estimates are presented. Rivalry between oppositely moving gratings allows eye movements to be recorded alongside active reports \citep{frassle2014rivalry}. A delayed response alone could reflect a sensory change, a memory effect or a reporting strategy. An earlier measure with an independent calibration can put some of these accounts into conflict.

\paragraph{Protocol.}
Calibrate a signed eye-movement decoder on a separate block of unambiguous physical alternations. On each test trial randomly cue attention to one grating or give a neutral cue. Record gaze over a fixed interval before any response key is assigned. Afterwards reveal a randomized key assignment and ask for the last dominant direction, allowing a mixed or undecided response. Interleave matched physical-replay trials with the same cues and key procedure. Give neutral trials an independently randomized target sign. Predeclare the timing, replay matching, uncertainty in calibration transfer, missing-trial handling, exclusions and analysis population.

\paragraph{One estimand for both rivals.}
Within each participant let $Y$ be the signed pre-key decoder score, positive toward the cued direction or the assigned neutral target. For mode $m$ (rivalry or replay), let $d_m$ be the mean cued score minus the mean neutral score. The target population contrast is
\[
 \Delta=\mathbb E_{\rm participants}\,[d_{\rm rivalry}-d_{\rm replay}] .
\]
Use the participant as the independent replication unit; repeated trials and gaze samples refine that participant's estimate. For a specified response-only account predeclare $|\Delta|\le\epsilon$. A specified sensory-selection account predicts $\Delta\ge\delta$ with $\delta>\epsilon$. Set both limits before examining rivalry outcomes and base them on the same contrast, including calibration uncertainty. If separate constituent effects are each bounded by $\epsilon_0$, the bound on their difference can be as large as $2\epsilon_0$.

\paragraph{Interpretation and rival explanations.}
An interval wholly above $\epsilon$ rejects the specified response-only bound. An interval wholly below $\delta$ rejects the specified sensory minimum. An estimate between the ranges may challenge both models. An interval spanning both ranges leaves them unresolved. Report the eye and button distributions, calibration errors, and sensitivity to missing trials. Cue-driven oculomotor attention, memory and failure of the replay calibration to transfer remain alternative explanations for a positive contrast. No outcome alone measures the presence of feeling. The design follows the logic of testing prespecified rival predictions on common data, also used in the Cogitate collaboration \citep{cogitate2025}.

\paragraph{Where the programme reaches.}
A physical identification requires convergence of independently prepared probes under controlled interventions. For example, the \source{network/README.md}{network reconstruction} obtains directed distances from minimal delivery costs with subadditive composition, and common limiting distances under calibrated agreement and refinement bounds. Reversible transports and specified local frame changes give a conditional gauge action. The results presuppose their devices and dynamics. They do not select a physical Standard Model, a four-dimensional spacetime or Einstein evolution from a seven-level register. Likewise, identifying an operational capability with experience remains an open bridge. A new measurement bears on that bridge when its predicted difference is specified independently of the outcome used to calibrate it.
